\documentclass[11pt]{article}
\usepackage{varnernotes}

\newcommand{\D}{\mathcal{D}}

\begin{document}

\lecturenotes{2a}{Pricing United States Treasury Securities}{Fall 2026}{September 1, 2026}{Jeffrey D. Varner}

\learningobjectives
  {Classify Treasury securities}
  {Distinguish Treasury bills, notes, and bonds using their maturities and promised cash-flow structures.}
  {Derive security prices}
  {Derive zero-coupon and coupon-bearing Treasury prices from the zero-NPV condition under a stated discounting convention.}
  {Evaluate risk claims}
  {Identify which risks the risk-free assumption excludes and which risks still matter.}

\section{Why Does the Treasury Borrow?}

When federal expenditures exceed tax and other receipts, the U.S. Department of the Treasury finances the resulting deficit by issuing debt. A Treasury security is therefore a contract: an investor supplies dollars today, and the federal government promises a specified pattern of future payments. The phrase \emph{fixed income} describes those contractual payments; it does not imply that the security's market price remains fixed.

Marketable Treasury securities differ mainly in maturity and promised payments (\tabref{treasury-classes}). Bills make one payment at maturity. Notes and bonds pay coupons and then repay principal. We must identify those cash flows before calculating a price.

\begin{table}[ht]
  \centering
  \caption{Current standard features of marketable nominal Treasury bills, notes, and bonds. The term is the time from issuance to maturity; market prices may change before maturity.}
  \label{tab:treasury-classes}
  \small\sffamily
  \renewcommand{\arraystretch}{1.2}
  \begin{tabular}{@{}p{0.16\textwidth}p{0.25\textwidth}p{0.34\textwidth}p{0.16\textwidth}@{}}
    \hline
    \textbf{Security} & \textbf{Standard term} & \textbf{Promised cash flows} & \textbf{Interest timing} \\
    \hline
    Treasury bill & 4--52 weeks & Face value at maturity & At maturity \\
    Treasury note & 2, 3, 5, 7, or 10 years & Semiannual coupons and face value & Every six months \\
    Treasury bond & 20 or 30 years & Semiannual coupons and face value & Every six months \\
    \hline
  \end{tabular}
\end{table}

Treasury auctions use bids to determine an issue rate, yield, or discount margin. Noncompetitive bidders accept the auction result. Competitive bidders state the rate or yield they will accept. The Treasury accepts competitive bids from the lowest financing cost upward until the offering is filled, and all successful bidders receive the same accepted rate, yield, or margin. We can value the security once those terms and the discounting convention are known.

\section{Cash-Flow Structure Comes First}

From the investor's perspective, purchasing a Treasury security means paying today and receiving money later. The timelines show the key difference: a bill makes one future payment, while a note or bond pays coupons and then repays principal (\figref{treasury-cash-flows}). An arrow leaving a timeline marker is money paid; an arrow entering it is money received.

\begin{figure}[ht]
  \centering
  \begin{tikzpicture}[x=1.45cm,y=1cm,>=Latex,font=\sffamily\small]
    \node[anchor=east,font=\bfseries] at (-0.15,1.45) {Bill};
    \draw[vnink,line width=0.8pt,-{Latex[length=3mm]}] (0,1.45)--(6.4,1.45) node[right]{time};
    \node[draw=vnink,circle,fill=white,minimum size=5mm,inner sep=0pt] (b0) at (0.45,1.45) {};
    \node[draw=vnink,circle,fill=white,minimum size=5mm,inner sep=0pt] (bT) at (5.8,1.45) {};
    \node[below=4pt] at (b0) {$t=0$};
    \node[below=4pt] at (bT) {$t=T$};
    \draw[vncarnelian,line width=1.3pt,-{Latex[length=3mm]}] (b0.north)--(0.45,2.35);
    \node[right,text=vncarnelian] at (0.45,1.92) {$-V_B$};
    \draw[vnlink,line width=1.3pt,-{Latex[length=3mm]}] (5.8,2.35)--(bT.north);
    \node[right,text=vnlink] at (5.8,1.92) {$+V_P$};

    \node[anchor=east,font=\bfseries] at (-0.15,-0.55) {Note/bond};
    \draw[vnink,line width=0.8pt,-{Latex[length=3mm]}] (0,-0.55)--(6.4,-0.55) node[right]{time};
    \foreach \x/\lab/\unit in {0.45/0/n0,1.55/1/n1,2.65/2/n2,4.65/N-1/nm,5.8/N/nN}{
      \node[draw=vnink,circle,fill=white,minimum size=5mm,inner sep=0pt] (\unit) at (\x,-0.55) {};
      \node[below=4pt] at (\unit) {$j=\lab$};
    }
    \node at (3.65,-0.53) {$\cdots$};
    \draw[vncarnelian,line width=1.3pt,-{Latex[length=3mm]}] (n0.north)--(0.45,0.35);
    \node[right,text=vncarnelian] at (0.45,-0.08) {$-V_B$};
    \foreach \unit/\x in {n1/1.55,n2/2.65,nm/4.65}{
      \draw[vnlink,line width=1.3pt,-{Latex[length=3mm]}] (\x,0.35)--(\unit.north);
      \node[right,text=vnlink] at (\x,-0.08) {$+C$};
    }
    \draw[vnlink,line width=1.3pt,-{Latex[length=3mm]}] (5.8,0.35)--(nN.north);
    \node[right,text=vnlink] at (5.8,-0.08) {$+C+V_P$};
  \end{tikzpicture}
  \caption{Cash flows from the investor's perspective. The purchase price $V_B$ is paid at the valuation date. A bill returns only its face value $V_P$ at maturity time $T$. A note or bond pays coupon amount $C$ at each coupon index $j=1,\ldots,N$ and adds the face value to the final coupon at index $N$.}
  \label{fig:treasury-cash-flows}
\end{figure}

The first timeline in \figref{treasury-cash-flows} describes a zero-coupon Treasury bill (\defnref{treasury-bill}). It separates the price paid today from the face value received at maturity.

\begin{definition}{Treasury bill}{treasury-bill}
Let time $0$ be the purchase date and let the maturity time be $T>0$ years. A Treasury bill is a zero-coupon debt security with purchase price $V_B$ dollars at time $0$ and face value $V_P$ dollars paid at time $T$. The term \emph{zero coupon} means that no contractual interest payment occurs between purchase and maturity. The nominal dollar gain at maturity is $V_P-V_B$; it is not itself a return. The continuously compounded holding-period growth rate and its associated dimensionless log return are
\begin{equation}
  \label{eq:bill-growth}
  g_B:=\frac{1}{T}\log\!\left(\frac{V_P}{V_B}\right),
  \qquad
  r_B:=Tg_B=\log\!\left(\frac{V_P}{V_B}\right).
\end{equation}
A quoted Treasury yield follows its own market convention and should not be identified with $g_B$ without the appropriate compounding and day-count conversion.
\end{definition}

The absence of intermediate cash flows makes a bill the cleanest application of L1b's rule \emph{convert, then add}. Under a specified model, its purchase price is the present value of the single promised receipt (\thmref{bill-price}).

\begin{theorem}{Zero-coupon price under periodic compounding}{bill-price}
Let the positive integer $n$ be the number of compounding intervals per year, let the maturity time be $T>0$ years, and let the constant nominal annual yield used for valuation be $y$ with $1+y/n>0$. Define its continuously compounded equivalent (with units of inverse years) by $g_y:=n\log(1+y/n)$ and the forward accumulation factor by
\begin{equation}
  \label{eq:yield-growth-conversion}
  \D_{T,0}(y;n):=\left(1+\frac{y}{n}\right)^{nT}=e^{g_yT}.
\end{equation}
The real exponent $nT$ need not be an integer for a zero-coupon claim. If the bill's face value $V_P$ is paid with certainty at maturity and taxes, transaction costs, and other market frictions are ignored, then the zero-NPV purchase price $V_B$ is given by \eqnref{bill-price}:
\begin{equation}
  \label{eq:bill-price}
  V_B=\D_{T,0}^{-1}(y;n)V_P
  =\left(1+\frac{y}{n}\right)^{-nT}V_P
  =e^{-g_yT}V_P.
\end{equation}
Here $\D_{T,0}^{-1}(y;n):=1/\D_{T,0}(y;n)$ converts maturity-date dollars to time-0 dollars. Under this flat periodic-yield model, \eqnref{bill-price} implies $g_B=g_y$; market quotation conventions may imply a different conversion.
\end{theorem}

The bill-pricing relation in \eqnref{bill-price} follows by setting the investor's NPV equal to zero and solving for the only time-0 cash flow:
\begin{align*}
  0
    &= -V_B+\D_{T,0}^{-1}(y;n)V_P,
    && \text{zero NPV at the purchase price},\\
  V_B
    &= \D_{T,0}^{-1}(y;n)V_P,
    && \text{isolate the time-0 price.}
\end{align*}

The compounding frequency $n$ and nominal yield $y$ in this course model must be stated explicitly. Treasury bill auction quotations also use market-specific bank-discount and investment-rate conventions; those quotation rules should not be silently identified with an arbitrary compound-interest model or with the continuously compounded growth rate $g_y$.

The companion notebook applies \thmref{bill-price} to auction records and makes the quotation convention visible in the calculation.

\notebooklink
  {L2a: Treasury Bill Pricing Worked Example}
  {https://github.com/varnerlab/CHEME-5660-CourseRepository-Fall-2026/blob/main/lectures/week-2/L2a/CHEME-5660-L2a-TBills-Pricing-Fall-2026-WorkedExample.ipynb}
  {Compute zero-coupon bill prices from auction inputs and compare model values with reported prices.}

\section{Coupon-Bearing Notes and Bonds}

Treasury notes and bonds spread value across time rather than concentrating it at maturity. That stream creates both an income component and greater sensitivity to how discount rates vary across maturities.

The repeated-payment structure in \figref{treasury-cash-flows} defines a coupon-bearing Treasury security (\defnref{coupon-security}) and distinguishes coupon income from principal repayment.

\begin{definition}{Coupon-bearing Treasury security}{coupon-security}
Let the positive integer $n$ be the number of coupon payments per year, let the maturity time be $T>0$ years with $N=nT$ coupon dates, let the face value be $V_P$ dollars, and let the annual coupon rate be $c$. A nominal coupon-bearing Treasury security pays the coupon amount $C:=(c/n)V_P$ dollars at every coupon date $j=1,\ldots,N$ and repays the face value $V_P$ with the final coupon. Treasury notes and bonds share this cash-flow structure but have different standard maturity ranges.
\end{definition}

Each promised receipt must be converted separately because it belongs to a different dated currency. Under a flat nominal yield, adding those time-0 values produces the coupon-security pricing relation (\thmref{coupon-price}).

\begin{theorem}{Coupon-security price under a flat yield}{coupon-price}
Let the positive integers $n$ and $N=nT$ denote the coupon frequency per year and total number of coupons over the maturity time $T$ years. Let $V_P$ be the face value, let $C$ be the coupon amount, and let the constant nominal annual yield used for valuation be $y$ with $1+y/n>0$. Define the accumulation factor to coupon date $j$ by $\D_{j,0}(y;n):=(1+y/n)^j$. If all promised payments are made with certainty and taxes, transaction costs, accrued-interest conventions, and other market frictions are ignored, then the zero-NPV purchase price is given by \eqnref{coupon-price}:
\begin{equation}
  \label{eq:coupon-price}
  V_B=\sum_{j=1}^{N}\D_{j,0}^{-1}(y;n)C+\D_{N,0}^{-1}(y;n)V_P.
\end{equation}
The summation index $j$ enumerates coupon dates, and each reciprocal factor converts that date's dollars to time-0 dollars.
\end{theorem}

The two terms in \eqnref{coupon-price} have simple meanings: the sum is the present value of the coupons, and the final term is the present value of principal. The zero-NPV derivation is therefore:
\begin{align*}
  0
    &= -V_B+\sum_{j=1}^{N}\D_{j,0}^{-1}(y;n)C
       +\D_{N,0}^{-1}(y;n)V_P,
    && \text{value every signed cash flow at time 0},\\
  V_B
    &= \sum_{j=1}^{N}\D_{j,0}^{-1}(y;n)C
       +\D_{N,0}^{-1}(y;n)V_P,
    && \text{solve for the purchase price.}
\end{align*}

The note-and-bond notebook evaluates this cash-flow sum using actual auction data; it is placed here because its code implements the coupon-bearing model directly.

\newpage
\thispagestyle{fancy}
\notebooklink
  {L2a: Treasury Note and Bond Pricing Worked Example}
  {https://github.com/varnerlab/CHEME-5660-CourseRepository-Fall-2026/blob/main/lectures/week-2/L2a/CHEME-5660-L2a-NotesBonds-Pricing-Fall-2026-WorkedExample.ipynb}
  {Construct coupon cash flows, discount each payment, and compare the resulting value with auction prices.}

\section{What Does ``Risk-Free'' Suppress?}

Calling a Treasury security risk-free is a modeling statement, not a description of every possible outcome. In a one-period valuation model, the phrase usually means that the promised dollar payment is treated as known and default is excluded. That assumption makes Treasury yields useful benchmark rates, but it does not conserve purchasing power or market value.

\textbf{Credit and political risk.}\enspace The probability that the U.S. Treasury fails to make a promised nominal payment is generally treated as very low, but institutional and political disruptions are not logical impossibilities.

\textbf{Interest-rate and liquidity risk.}\enspace A holder who sells before maturity receives the prevailing market price, not the contractual face value. When market discount rates rise, the present value of fixed promised payments falls; L2b develops this sensitivity precisely. Treasury markets are highly liquid, but bid-ask spreads and market depth can still affect an immediate sale price.

\textbf{Inflation and reinvestment risk.}\enspace Nominal repayment does not guarantee purchasing power. Coupon-bearing securities also require the investor to reinvest intermediate coupons at rates that are unknown when the security is purchased.

Always qualify the phrase \emph{risk-free}: which risk is excluded, over what holding period, and under which assumptions?

\keytakeaways
  {In this lecture, we represented Treasury bills, notes, and bonds as abstract assets and derived their prices by discounting promised cash flows to a common valuation date.}
  {Structure determines value}
  {Bills contain one maturity receipt, while notes and bonds contain coupon streams plus principal; the cash-flow structure must be identified before pricing.}
  {Price means zero NPV}
  {Under stated assumptions, a Treasury's purchase price is the amount that makes the investor's discounted signed cash flows sum to zero.}
  {Risk-free is conditional}
  {Treating nominal payments as certain suppresses default risk but does not eliminate interest-rate, inflation, reinvestment, or liquidity concerns.}
  {Next, use these pricing relations to determine how a bond's value changes when its yield, coupon, maturity, or payment timing changes.}

\begin{readinglist}
  \item U.S. Treasury, \href{https://www.treasurydirect.gov/marketable-securities/treasury-bills/}{\emph{Treasury Bills}}, for current terms, cash-flow conventions, and auction links.
  \item U.S. Treasury, \href{https://www.treasurydirect.gov/marketable-securities/treasury-notes/}{\emph{Treasury Notes}} and \href{https://www.treasurydirect.gov/marketable-securities/treasury-bonds/}{\emph{Treasury Bonds}}, for current maturities and coupon-payment conventions.
  \item U.S. Treasury, \href{https://www.treasurydirect.gov/auctions/how-auctions-work/}{\emph{How Auctions Work}}, for competitive and noncompetitive bidding mechanics.
\end{readinglist}

\vfill
{\sffamily\footnotesize\color{vnmuted}\textbf{Educational use and risk notice.} These notes are for instruction only and do not constitute investment advice. Financial decisions involve risk, and model outputs depend on assumptions.}

\end{document}
